% Generated draft: requires mathematical and source review.
\documentclass[11pt]{article}
\usepackage[margin=25mm]{geometry}
\usepackage{fontspec}
\setmainfont{Latin Modern Roman}
\usepackage{amsmath,amssymb,mathtools,unicode-math}
\setmathfont{Latin Modern Math}
\usepackage{enumitem,needspace,xurl,hyperref}
\setlength{\parindent}{0pt}
\setlength{\parskip}{5pt}
\setlength{\emergencystretch}{3em}
\newcommand{\N}{\mathbb{N}}
\newcommand{\Z}{\mathbb{Z}}
\newcommand{\Q}{\mathbb{Q}}
\newcommand{\R}{\mathbb{R}}
\newcommand{\C}{\mathbb{C}}
\newcommand{\F}{\mathbb{F}}
\newcommand{\A}{\mathbb{A}}
\newcommand{\PP}{\mathbb{P}}
\newcommand{\E}{\mathbb{E}}
\newcommand{\eps}{\varepsilon}
\newcommand{\dd}{\,\mathrm{d}}
\begin{document}
\section*{UnsolvedMath ID: 701741}
\textbf{Source problem code:} catalogue entry
\subsection{Definitions and mathematical statement}
Let \(\mathfrak{X}\) be a class of finite groups closed under taking subgroups, homomorphic images, and extensions. Closure under extensions means that a finite group \(E\) belongs to \(\mathfrak{X}\) whenever it has a normal subgroup \(D\) with \(D\in\mathfrak{X}\) and \(E/D\in\mathfrak{X}\).

An \(\mathfrak{X}\)-group is a group belonging to \(\mathfrak{X}\). A subgroup \(M\leq L\) is \(\mathfrak{X}\)-maximal in \(L\) if \(M\in\mathfrak{X}\) and there is no subgroup \(U\leq L\) such that \(M<U\) and \(U\in\mathfrak{X}\). In particular, this does not require \(M\) to be a maximal proper subgroup of \(L\).

A subgroup \(K\leq L\) is subnormal if there is a finite chain
\[
K=K_0\trianglelefteq K_1\trianglelefteq\cdots\trianglelefteq K_r=L.
\]
For a finite group \(K\), a subgroup \(A\leq K\) is \(\mathfrak{X}\)-submaximal if there exist an injective homomorphism \(\iota:K\to K^*\) and an \(\mathfrak{X}\)-maximal subgroup \(M\leq K^*\) such that \(\iota(K)\) is subnormal in \(K^*\) and
\[
\iota(A)=\iota(K)\cap M.
\]
The supplied definition calls \(K^*\) a group without explicitly requiring it to be finite. This draft retains that wording; whether the intended definition restricts \(K^*\) to finite groups requires confirmation.

\textbf{Question.} For every such class \(\mathfrak{X}\), every finite group \(G\), and every characteristic subgroup \(N\) of \(G\), suppose that all \(\mathfrak{X}\)-submaximal subgroups of \(N\) are conjugate in \(N\): for any two such subgroups \(A,B\leq N\), there exists \(n\in N\) such that \(B=n^{-1}An\). Here characteristic means that \(\alpha(N)=N\) for every automorphism \(\alpha\) of \(G\). Must it follow that, for every \(\mathfrak{X}\)-submaximal subgroup \(H\leq G\), the subgroup
\[
HN/N=\{hN:h\in H\}\leq G/N
\]
is \(\mathfrak{X}\)-submaximal in \(G/N\)?

\textbf{Supplied qualifications.} The source states that replacing the characteristic hypothesis by normality does not suffice, even if \(N\in\mathfrak{X}\), or if \(N\) contains no nontrivial \(\mathfrak{X}\)-subgroup. The latter condition means that \(U\leq N\) and \(U\in\mathfrak{X}\) imply \(U=1\). The source also reports a positive answer when \(N\) is the \(\mathfrak{F}\)-radical of \(G\) for a Fitting class \(\mathfrak{F}\). Here a Fitting class is a nonempty class of finite groups closed under normal subgroups and products of normal subgroups belonging to the class; its radical in \(G\) is the largest normal subgroup of \(G\) belonging to that class. These counterexample and positive-result assertions are supplied background, not independently verified results in this draft.
\subsection{Short English statement}
Suppose all \(\mathfrak{X}\)-submaximal subgroups of a characteristic subgroup \(N\) of a finite group \(G\) are conjugate within \(N\). Must the image of every \(\mathfrak{X}\)-submaximal subgroup of \(G\) be \(\mathfrak{X}\)-submaximal in \(G/N\)? The supplied record labels this problem open. The intended finiteness requirement on the ambient group in the definition needs confirmation.
\subsection{Sources}
\begin{itemize}
\item Local UnsolvedMath registry, permanent ID \(701741\), exact source problem code \texttt{catalogue entry}, title \emph{Kourovka Notebook Problem 20.83}. The supplied record labels the problem \emph{open}; that status is not independently verified here. The record gives creation and update timestamps \texttt{2026-10-16T14:00:00Z}.
\item \emph{The Kourovka Notebook (Unsolved Problems in Group Theory)}, Problem 20.83, 20th issue (2022), also identified in the supplied background with the 21st issue edition (2026), \texttt{arXiv:1401.0300v48}. The supplied collection description names E. I. Khukhro and V. D. Mazurov as editors.
\item Proposed by D. O. Revin and A. V. Zavarnitsine; proposed year supplied as 2022.
\item The normal-subgroup counterexamples and the positive result for a Fitting-class radical are reported in the supplied statement without separate supporting references. No external source was consulted for this draft.
\end{itemize}
\end{document}
